Consider an Alloy model M, with a command C, of the form run/check P for scope s. When fed into the Analyzer, there are 4 possible outcomes:

\begin{itemize}
  \item Instance found - this occurs when C is a run command, and the constraints of M and P are satisfied by an instance I, which is shown by the Analyzer.
  \item No Instances found - this occurs when when C is a run command, and there is no instance I capable of satisfying the constraints of M and P. This includes the case where the constraints of just M are unsatisfiable. 
  \item Counterexample found - this occurs when C is a check command, and there is an instance I which satisfies the constraints of M, and violates the constraints of P.
  \item No Counterexample found - this occurs when C is a check command, and every instance I which satisfies the constraints of M, also satisfies the constraints of P. This includes the case where the constraints of just M are unsatisfiable. 
\end{itemize}

A check command in Alloy can be turned into an equivalent run command by negating the body of the command. $check p for s$ is equivalent to $run ~p for s$. Going forward, the command is assumed to be a run command.

The goal of a translation to TLA+ is that when the translated model is run using TLC, the output produced by TLC can be used to infer the output the Analyzer would have produced, had the original model been run on the Analyzer. A translation scheme needs to map from the output of TLC to the output of the Analyzer.

A translation is correct, if and only if the Analyzer's outputs correspond to the outputs produced by TLC. A translation is efficient, if the time taken to run a translated model using TLC is lower than the time taken to run the original model using the Analyzer.